% chapters/04_anleitung.tex

\chapter{Anleitung zum Weiterarbeiten}
\label{chap:anleitung}

Die Befehle entsprechen \doc{docs/quickstart.md}; dort steht zu jedem Schritt, was danach zu
sehen sein muss. Alle Wege lassen sich auch über \texttt{./run.sh} starten.

\section{Zugänge}
\label{sec:zugaenge}

\begin{itemize}
  \item Hugging-Face-Token (Pflicht): Lizenz von \texttt{nvidia/GR00T-N1.6-3B} akzeptieren; für
    N1.7 zusätzlich Zugang zu \texttt{nvidia/Cosmos-Reason2-2B}.
  \item \ac{wandb}-Schlüssel: lokal optional, auf KISSKI Pflicht (sonst
    \texttt{USE\_WANDB=0}).
  \item KISSKI-Konto für das Training, ein Rechner mit RT-Core-GPU für die Simulation
    (\cref{tab:rechenumgebungen}).
  \item Namensräume: Docker Hub \texttt{lucam03} (Images), GitHub \texttt{lucam06}
    (GR00T-Fork), Hugging Face \texttt{luca-mue} (Checkpoints).
\end{itemize}

\section{Training}
\label{sec:anl-training}

\begin{minted}{bash}
# lokal: zieht das Image, startet Container groot-train (Standard 30 000 Schritte)
export HF_TOKEN=hf_... WANDB_API_KEY=...
./Training/setup_and_train_dockerhub_pull.sh        # N1.7: GROOT_VERSION=1.7 davor
./Training/setup_and_train_dockerhub_pull.sh --resume  # nach Stopp fortsetzen

# KISSKI (Login-Knoten), einmalig das Image holen:
module load apptainer
apptainer pull $HOME/images/projekt-humanoider-roboter.sif \
    docker://lucam03/projekt-humanoider-roboter:latest
export HF_TOKEN=hf_... WANDB_API_KEY=... GLOBAL_BATCH_SIZE=32
sbatch --export=ALL Training/kisski_submit.sh        # 4x A100, 44 000 Schritte
\end{minted}

Varianten werden über Umgebungsvariablen gewählt: \texttt{TUNE\_VISUAL=1} trainiert den
Bild-Encoder mit, \texttt{USE\_COTRAIN=1} mischt gerenderte Daten bei,
\texttt{TRAIN\_TEST\_SPLIT=1} hält die Episoden 240--300 zurück. \texttt{Training/kisski\_chain.sh}
reiht mehrere Läufe mit anschließender Checkpoint-Auswertung hintereinander ein. Deren
Standard sind 20 Testepisoden in voller Länge (\texttt{EVAL\_NUM\_TRAJ=20
EVAL\_STEPS=750}); das frühere kurze Protokoll mit 6 Episoden à 300 Schritten lieferte
Scheinoptima (\cref{sec:erg-training}). Der fertige Checkpoint wird ohne Optimizer hochgeladen:

\begin{minted}{bash}
python Simulation/scripts/upload_checkpoint.py \
    --checkpoint <...>/checkpoint-44000 --repo luca-mue/groot-g1dex3-<lauf>-ckpt44000
\end{minted}

\section{Checkpoint messen}
\label{sec:anl-messen}

Auf dem Sim-Server erst prüfen, ob Image und GPU passen, dann Gate und Closed Loop:

\begin{minted}{bash}
./Simulation/server_rl_run.sh preflight
# visuelles Gate, Seeds 0/1/2 (Container, GROOT_VERSION, GPU, HF-Repo ohne Konto)
tools/sim_visgate_seeds.sh groot-rl 1.6 0 groot-g1dex3-cotrain-v3-cube4-ckpt44000
python3 tools/sim_visgate_seeds_summary.py ~/groot-rl-data/sim_results \
    cotrain-v3-cube4-ckpt44000
# Closed Loop nach Protokoll (20 Testlagen x 120 s, Szene aus Lauf 70)
SIM_CUBE_EDGE=0.04 tools/sim_cl_watchdog.sh cl_meinlauf groot-rl 0 1.6 \
    luca-mue/groot-g1dex3-cotrain-v3-cube4-ckpt44000 \
    /data/checkpoints/cotrain-v3-cube4-ckpt44000
python3 tools/sim_cl_summary.py ~/groot-rl-data/sim_results/results_<...>.json
\end{minted}

\texttt{SIM\_CUBE\_EDGE=0.04} gilt nur für Modelle, die mit 4-cm-Würfeln co-trainiert wurden.
Ein Treffer bei \texttt{lifted\_clear} zählt erst, wenn das Video unter
\texttt{\textasciitilde/groot-rl-data/sim\_videos/<tag>/} einen Griff zeigt
(\cref{sec:grenzen-metriken}). Die Prüfbilder des Gates liegen nur auf dem IKR-Server; auf
einem neuen Server müssen sie mit \texttt{server\_rl\_run.sh render} neu erzeugt werden,
das dy3-Layout mit \doc{tools/layout_shift.py --dy 0.03} (Befehle in \doc{docs/quickstart.md}).

\section{Häufige Fehler}
\label{sec:anl-fehler}

\begin{table}[htb]
  \centering
  \begin{tabularx}{\linewidth}{XX}
    \toprule
    Symptom & Ursache und Abhilfe \\
    \midrule
    KISSKI-Job läuft ohne Token & \texttt{--export=ALL} fehlt; der Login-Knoten setzt
      \texttt{SBATCH\_EXPORT=none} \\
    KISSKI-Job bricht mit \enquote{invalid device ordinal} ab & vier GPUs auf zwei Knoten
      verteilt; \texttt{\#SBATCH -N 1} setzen (Lauf~75) \\
    Checkpoint lädt nicht, Warnung zu \texttt{model\_type} & falsche \texttt{GROOT\_VERSION};
      N1.6 und N1.7 sind nicht gegenseitig ladbar \\
    Training bricht mit Speicherfehler ab & \texttt{GLOBAL\_BATCH\_SIZE} halbieren; bei
      Arbeitsspeicher-Abbruch weniger Datenlade-Prozesse je GPU \\
    Closed Loop hängt ohne Fehlermeldung & bekannter Hänger (Läufe 89, 90, 97);
      \texttt{sim\_cl\_watchdog.sh} startet nach 15~min neu \\
    \bottomrule
  \end{tabularx}
  \caption{Die fünf häufigsten Fehler. Weitere: \doc{docs/fehlerbehebung.md}.}
  \label{tab:anl-fehler}
\end{table}
